\section{The minor-arc bound}\label{sec:minor}

This section proves the minor-arc half of the dichotomy. The write-up
forced a genuine correction to the working notes, recorded in
Remark~\ref{rem:M2-strengthening}: the planned hypotheses (a single
convergent denominator near $R$) do not suffice, and the planned
conclusion (a polylogarithmic bound for every minor $p$) is not what the
correct proof delivers. The true shape is a \emph{Diophantine-functional}
bound: $|V_p|$ is controlled by the convergent-gap structure of $q/p$
below scale $R$, polylogarithmic for typical pairs and
$O(\sqrt R\,\Llog^{O(1)})$ --- still far within the budget of
Proposition~\ref{prop:digit-assembly} --- on the whole minor set.

Throughout: $\alpha:=q/p\in(0,1)$ in lowest terms,
$\beta:=\fr{\lambda p/q}$, $\theta_\mu=\mu\alpha-\beta\ (\mathrm{mod}\ 1)$,
$R=x/(pq)$, $\Llog=\log x$, $\langle\mu\rangle=\min(\mu,p-\mu)$, and
$B\ge1$ a level parameter. Denominators of continued-fraction convergents
of $\alpha$ are denoted $s$, with successor $s^{+}$ (for the rational
$\alpha=q/p$ the final denominator is $p$; everything below concerns
$s<p$).

\begin{definition}[gap functional; minor conditions]\label{def:minor-conditions}
Let $s_0=s_0(\alpha)$ be the largest convergent denominator $\le R$, and
\[
\mathcal D_R(\alpha)\;:=\;
\sum_{\substack{s\ \mathrm{convergent\ denominator}\\ s\le R}}
\sqrt{\frac{s^{+}}{s}}\,.
\]
Call the band prime $p$ \emph{minor for $(q,\lambda)$ at level $B$} if:
\begin{itemize}
\item[(M1)] $\displaystyle
\frac{1}{\langle\mu\rangle}\,
\min\Big(R,\ \frac{1}{2\nm{\theta_\mu}}\Big)\ \le\ \Llog^{B}$
\quad for every $\mu$ with $0<\langle\mu\rangle\le R$;
\item[(M2$^*$)] every convergent denominator $s\le R$ of $\alpha$
satisfies $s^{+}/s\le R\,\Llog^{-B}$.
\end{itemize}
\end{definition}

\begin{remark}
(M2$^*$) is \emph{not} the smallness of $\mathcal D_R$: the first
partial quotient of $\alpha=q/p$ is $\asymp p/q$, so for $u>u'$
\emph{every} pair has the structural gap $s=1$, $s^{+}\asymp p/q$, and
$\mathcal D_R\ge\sqrt{p/q}$ wholesale. The threshold $R\Llog^{-B}$ is
calibrated so that this structural gap is automatically exempt:
$p/q\le R\Llog^{-B}$ amounts to $x^{2u-1}\le\Llog^{-B}$, which holds on
the band after the top trim $u\le\tfrac12-\eta$. Gap ratios beyond
$R\Llog^{-B}$ are genuinely rare (Proposition~\ref{prop:familyB}).
\end{remark}

\begin{lemma}[continued-fraction toolkit]\label{lem:cf}
\statusP\quad
Let $s\le M<s^{+}$ for a convergent denominator $s$ of $\alpha$, with
convergent $a/s$. Then:
\begin{enumerate}
\item (best approximation)
$\min_{0<k<s^{+}}\nm{k\alpha}=\nm{s\alpha}\in
\big[\tfrac{1}{s+s^{+}},\tfrac1{s^{+}}\big]$;
\item (lattice proximity) for any $\mu_0$ and $0\le k<s$,
\[
\theta_{\mu_0+k}\;=\;\theta_{\mu_0}+\frac{ka}{s}+O^*\Big(\frac1{s^{+}}\Big)
\pmod 1,
\]
where $|O^*(t)|\le t$; the $s$ sites $\theta_{\mu_0}+ka/s$ are distinct
points of a translated $\frac1s$-grid;
\item (run shells) consequently, within any window of $s$ consecutive
$\mu$, the $k$-th smallest value of $\nm{\theta_\mu}$ is
$\ge\frac{k-4}{2s}$ for $k\ge5$;
\item (across runs) writing a window of $\le M$ consecutive $\mu$ as
$\le M/s+1$ runs of $s$ consecutive $\mu$, the per-run nearest values
$\theta_{\mu_0+js}$ ($j=0,1,\ldots$) move along an arithmetic progression
of step $\eta:=\nm{s\alpha}\le1/s^{+}$, up to the run-internal
$O^*(1/s^{+})$.
\end{enumerate}
\end{lemma}

\begin{proof}
(1) is the best-approximation law together with the standard estimate
for $\nm{s\alpha}$ \cite[Ch.~I--II]{Khin}. (2): $k\alpha=ka/s+k\delta$
with $\delta=\alpha-a/s$, $|\delta|\le1/(ss^{+})$, $k<s$; distinctness
since $(a,s)=1$. (3): the $k$-th nearest \emph{site} of a
$\frac1s$-grid is at distance $\ge\fl{(k-1)/2}/s$, and the perturbation
$1/s^{+}$ can displace each point by at most one half grid-step in rank
terms; the stated bound is a safe form. (4): $\theta_{\mu_0+js}
=\theta_{\mu_0}+js\alpha=\theta_{\mu_0}\pm j\eta\bmod 1$.
\end{proof}

\begin{theorem}[D$^\dagger$-minor, functional form]\label{thm:minor}
\statusP\quad
Assume (M1) at level $B$, and $\sigma\equiv1$ (anatomy off). Then
\[
|V_p(\lambda)|\ \ll\ \Llog^{\,B+3}\big(1+\mathcal D_R(q/p)\big)
\;+\;\Llog^{2}\,\frac{R}{s_0(q/p)}\,.
\]
\end{theorem}

\begin{corollary}[minor bound]\label{cor:minor}
\statusP\quad
If $p$ is minor at level $B$ (i.e.\ also (M2$^*$)), then
$\mathcal D_R\le\Llog\sqrt{R\Llog^{-B}}$ and
$R/s_0\le s_0^{+}/s_0\le R\Llog^{-B}$, so
\[
|V_p(\lambda)|\ \ll\ \sqrt R\;\Llog^{\,B/2+4}\;+\;R\,\Llog^{\,2-B}.
\]
With $B\ge c+4$ and the corner trim $\sqrt R\ge\Llog^{B/2+c+5}$, this is
$\le R\Llog^{-c-1}$: within the budget of
Proposition~\ref{prop:digit-assembly}, with room.
\end{corollary}

\begin{proof}[Proof of Theorem~\ref{thm:minor}]
By Proposition~\ref{prop:two-freq} and
$|\hat f_\mu|\le\tfrac1{2\langle\mu\rangle}$,
\begin{equation}\label{eq:minor-start}
|V_p(\lambda)|\;\le\;\frac12\sum_{0<\mu<p}
\frac{1}{\langle\mu\rangle}\,X_\mu,
\qquad
X_\mu:=\min\Big(R,\ \frac{1}{2\nm{\theta_\mu}}\Big).
\end{equation}
Split $\{0<\mu<p\}$ into \emph{Region I} ($\langle\mu\rangle\le R$) and
\emph{Region II} ($\langle\mu\rangle>R$), and each region into dyadic
blocks: the $\mu$ with $\langle\mu\rangle\in(M,2M]$ form two intervals of
at most $M$ consecutive integers, with weight
$\langle\mu\rangle^{-1}\le1/M$. Fix such an interval $I$ at scale $M$ and
let $s$ be the largest convergent denominator $\le\min(M,R)$, so that
$s\le M<s^{+}$ in Region I.

\medskip\noindent\emph{Region I, within-run shells.} Partition $I$ into
$\le M/s+1$ runs of at most $s$ consecutive $\mu$. In each run,
Lemma~\ref{lem:cf}(3) gives, for the terms beyond the four nearest,
\[
\sum_{k\ge5}X_{\mu(k)}\;\le\;\sum_{5\le k\le s}\frac{s}{k-4}
\;\le\;2s\log x ,
\]
so the within-run shell totals contribute, with weight and run count,
\[
\frac1M\Big(\frac Ms+1\Big)\cdot 2s\log x\;\le\;4\log x .
\]

\medskip\noindent\emph{Region I, the near terms, two bounds.} It remains
to treat the $\le4$ nearest terms of each run --- at most $4(M/s+1)$
terms in the block. By (M1), each satisfies
$X_\mu/\langle\mu\rangle\le\Llog^{B}$, whence the \emph{floor bound}
\begin{equation}\label{eq:floor-branch}
\frac1M\sum_{\text{near terms}}X_\mu
\;\le\;4\Big(\frac Ms+1\Big)\Llog^{B}\Big/1
\quad\text{(each weighted term $\le\Llog^B$),}
\end{equation}
i.e.\ block near-term contribution $\le4(M/s+1)\Llog^{B}$. Alternatively,
by Lemma~\ref{lem:cf}(4) the near terms ride $\le4$ interleaved
arithmetic progressions of step $\eta\ge1/(2s^{+})$
(Lemma~\ref{lem:cf}(1)): in each progression at most two terms sit within
$\eta$ of the minimum --- these are again (M1)-floored, weighted
$\le\Llog^B$ each --- and the rest telescope:
\begin{equation}\label{eq:ap-branch}
\frac1M\sum_{\text{near terms}}X_\mu
\;\le\;8\Llog^{B}
+\frac1M\cdot8\,\frac{1}{\eta}\log x
\;\le\;8\Llog^{B}+\frac{16\,s^{+}}{M}\log x .
\end{equation}
Taking the minimum of \eqref{eq:floor-branch} and \eqref{eq:ap-branch}
and summing over the dyadic $M$ in the convergent gap $[s,\min(s^{+},R))$
--- the two branches cross at $M^*=\sqrt{s\,s^{+}\log x\,\Llog^{-B}}$,
with value $\ll\sqrt{(s^{+}/s)\,\Llog^{B}\log x}$, and decay
geometrically on either side ---
\[
\sum_{M\ \mathrm{dyadic}\ \in[s,\,s^{+})}
\min\big(\eqref{eq:floor-branch},\eqref{eq:ap-branch}\big)
\;\ll\;\sqrt{\frac{s^{+}}{s}}\;\Llog^{\,B/2+1}
\;+\;\Llog^{B+1} .
\]
Summing over the convergent gaps with $s\le R$ (at most $O(\log x)$ of
them) gives
\[
\text{Region I}\;\ll\;\Llog^{\,B/2+1}\,\mathcal D_R(\alpha)
+\Llog^{\,B+2}
\;\ll\;\Llog^{\,B+3}\big(1+\mathcal D_R(\alpha)\big)/\Llog .
\]

\medskip\noindent\emph{Region II.} Here $\langle\mu\rangle>R$, the weight
is $<1/R\le1/M\cdot(M/R)$, and (M1) is unavailable, but the trivial cap
$X_\mu\le R$ together with the weight makes every term $\le R/M\le1$.
Take $s=s_0$ (largest convergent $\le R$) and partition the block into
$\le M/s_0+1$ runs of $s_0$ consecutive $\mu$. Within-run shells as
before contribute $\le4\log x$ per block. The $\le4$ near terms per run
contribute, using $X_\mu\le R$ and the weight $1/M$,
\[
\frac1M\cdot4\Big(\frac M{s_0}+1\Big)R
\;=\;\frac{4R}{s_0}+\frac{4R}{M}
\;\le\;\frac{4R}{s_0}+4 .
\]
Summing the $O(\log x)$ dyadic blocks of Region II:
$\;\ll\;\log x\,\big(\log x+R/s_0\big)
\ll\Llog^2\big(1+R/s_0\big)$.

\medskip Combining the two regions in \eqref{eq:minor-start} yields the
theorem.
\end{proof}

\begin{remark}[what changed, and why it matters; a WP4 correction]
\label{rem:M2-strengthening}
The working notes (WP2.1/2.2) stated the minor lemma as: (M1) plus
``some convergent denominator in $[R\Llog^{-B},R\Llog^{B}]$'' implies
$|V_p|\ll\Llog^{B}\log x$. Carrying out the proof shows both halves need
repair. (i) The single-scale condition controls only the blocks at scale
$\approx R$; at scale $M\ll R$ the relevant structure is the convergent
gap \emph{at scale $M$}, and a long gap there lets a block carry mass
$\asymp\sqrt{(s^{+}/s)}$ in perfect consistency with (M1) --- the
``graze'' configuration: a run-progression approaching the origin at the
(M1)-floor. Hence the all-scale character of (M2$^*$) and of the
functional $\mathcal D_R$. (ii) The polylog conclusion is genuinely false
under (M1) plus any reasonable sparsity condition: the structural $s=1$
gap forces $\mathcal D_R\ge\sqrt{p/q}$ for the entire half-band $u>u'$.
What is true --- and suffices, with room, for
Estimate~D$^\dagger$ --- is the functional bound of
Theorem~\ref{thm:minor}, whose worst minor-set value is
$\sqrt R\,\Llog^{O(1)}$ and whose typical value is polylogarithmic
(the band average of $\mathcal D_R$ is $O(\Llog^2)$;
Proposition~\ref{prop:familyB}). The empirical record matches this
precisely: the observed minor-arc medians sit at fixed multiples of
$\sqrt R$-normalized $z\approx0.08$, i.e.\ square-root scale with small
constants, not at $\Llog/\!\sqrt R$ scale
(Appendix~\ref{app:numerics}). The WP2.2 conclusion ---
\emph{D$^\dagger$-digit closes modulo the two average-equidistribution
steps and the trims} --- survives with the repaired statements; the
intermediate claims did not.
\end{remark}
